\documentclass[11pt,a4paper]{article}
\usepackage[margin=1in]{geometry}
\usepackage[T1]{fontenc}
\usepackage{lmodern}
\usepackage{amsmath,amssymb,amsthm,mathtools,booktabs,microtype}
\usepackage{xurl}
\usepackage[hidelinks]{hyperref}
\hypersetup{pdftitle={A word model for three Schreier multiset recurrences},
  pdfauthor={Devin O'Keefe}}
\newtheorem{theorem}{Theorem}[section]
\newtheorem{proposition}[theorem]{Proposition}
\newtheorem{corollary}[theorem]{Corollary}
\newtheorem{lemma}[theorem]{Lemma}
\theoremstyle{remark}
\newtheorem{remark}[theorem]{Remark}
\newcommand{\A}{\mathcal A}
\newcommand{\Z}{\mathbb Z}
\newcommand{\Q}{\mathbb Q}
\newcommand{\ms}{\mathrm{ms}}
\DeclarePairedDelimiter{\floor}{\lfloor}{\rfloor}
\DeclarePairedDelimiter{\ceil}{\lceil}{\rceil}
\title{A word model for three Schreier multiset recurrences}
\author{Devin O'Keefe\\\normalsize Independent Researcher}
\date{}

\begin{document}
\maketitle
\begin{abstract}
We study uncoloured multisets $F$ of positive integers with maximum $n$
occurring once, each smaller element occurring at most twice, and
$|F|+r\le q\min F$, where $q\ge1$ and $0\le r<q$ are integers.
An encoding by words over $\{0,1,2\}$ gives a finite system of equations
for the generating functions indexed by $r$. For $r=0$ and $q=2,3,4$,
nine polynomial identities yield proofs of three conjectured recurrences,
together with their initial values and ranges of validity. We also prove
that the minimal orders of eventual homogeneous constant-coefficient
recurrences over $\Q$ are $3,5,5$, respectively. The encoding gives a
correspondence with restricted integer compositions, including compositions
of $3n+2$ into parts $3,4,5$ when $q=3$ and $r=0$.
\end{abstract}

\section{The families and the three recurrences}

For an integer $n\ge1$, let
\[
 U_n=\{1,1,2,2,\ldots,n-1,n-1,n\}.
\]
All multisets in this paper are uncoloured: the two available copies of a
lower label are indistinguishable. We write $F\le_{\ms}U_n$ for
multiset containment, allowing equality, and $|F|$ for the total number
of occurrences. For integers $q\ge1$ and $0\le r<q$, define
\begin{equation}\label{eq:family}
 \A_{q,r}(n)=\{F\le_{\ms}U_n:n\in F,\ |F|+r\le q\min F\},
 \qquad a_{q,r}(n)=|\A_{q,r}(n)|.
\end{equation}
Every member is nonempty and has a unique occurrence of its maximum $n$.
Set $\A_q(n)=\A_{q,0}(n)$ and $a_q(n)=a_{q,0}(n)$.

Chu et al.\ \cite[p.~19]{CGKMTV} conjectured recurrence relations for the
sequences $a_q(n)$ with $q=2,3,4$. We prove these recurrences using a
common word encoding.

\begin{theorem}\label{thm:main}
The sequences defined by \eqref{eq:family} have the following initial values:
\[
\begin{array}{c|c|c}
q&\text{initial values, starting at }n=1&\text{recurrence range}\\ \hline
2&1,2,4&n\ge4\\
3&1,3,6,13,31&n\ge6\\
4&1,3,8,18,41&n\ge6.
\end{array}
\]
On the indicated ranges, respectively,
\begin{align}
 a_2(n)&=a_2(n-1)+2a_2(n-2)+a_2(n-3),\label{eq:r2}\\
 a_3(n)&=3a_3(n-1)-3a_3(n-2)+4a_3(n-3)
          -2a_3(n-4)+a_3(n-5),\label{eq:r3}\\
 a_4(n)&=3a_4(n-1)-3a_4(n-2)+3a_4(n-3)
          +2a_4(n-4)+a_4(n-5).\label{eq:r4}
\end{align}
\end{theorem}

Specialising the counting formula of Chu et al.\
\cite[Section~5]{CGKMTV} to multiplicity at most two gives, for $q\ge2$,
\begin{equation}\label{eq:occupancy}
 a_q(n)=\sum_{K=1}^{2n-1}[z^{K-1}](1+z+z^2)^{\,n-\ceil{K/q}}.
\end{equation}
Indeed, at fixed cardinality $K$, all occupied labels are at least
$\ceil{K/q}$; the positions from there to $n-1$ have multiplicities $0,1,2$,
and one copy of $n$ is inserted. The exponent is nonnegative because
$K\le2n-1$ and $q\ge2$. At $n=1$, the convention $(1+z+z^2)^0=1$ and
$[z^0]1=1$ handles the zero-position case directly.

The sequence in \eqref{eq:r2} is already recorded as A141015, with an
additional initial zero, and its generating function is known
\cite{OEIS141015}. The parent composition sequence used below for $q=3$ is
A017818 \cite{OEIS017818}. Polynomial divisibility and extraction of periodic
subsequences are established tools; see, in this setting,
\cite[Section~3.2, especially Lemma~5]{CV}. Our contribution is a common
bijection between these multisets and words, including margins and the
maps relating the three constructions. The residue equations of this
model give short polynomial proofs of the three recurrences.

\section{A word encoding with a margin}

A word means a finite ordered list $w=(d_1,\ldots,d_L)$ with
$d_t\in\{0,1,2\}$; the empty word is allowed. Write
\begin{equation}\label{eq:cost}
 S(w)=\sum_{t=1}^L d_t,\qquad
 c_{q,r}(w)=L+\floor{\frac{S(w)+r}{q}}.
\end{equation}
Each position records a label, whether or not that label occurs in the multiset.

\begin{theorem}[Margin--word correspondence]\label{thm:encoding}
For every $q\ge1$, $0\le r<q$, and $n\ge1$, there is a bijection
\[
 E_{q,r,n}:\A_{q,r}(n)\longrightarrow
 \{w:c_{q,r}(w)=n-1\}.
\]
For $F$ of cardinality $K$, put $j=\floor{(K-1+r)/q}$. Its image is the
word of multiplicities of $F$ at the consecutive labels $j+1,\ldots,n-1$.
Conversely, a word of length $L$ and sum $S$ is decoded by putting
$j=\floor{(S+r)/q}$, inserting $d_t$ copies of $j+t$ for $1\le t\le L$,
and inserting one copy of $L+j+1$.
\end{theorem}

\begin{proof}
For positive $K$ and any integer label $x$,
\begin{equation}\label{eq:threshold}
 K+r\le qx\quad\Longleftrightarrow\quad
 K-1+r<qx\quad\Longleftrightarrow\quad
 \floor{\frac{K-1+r}{q}}<x.
\end{equation}
Thus every member of $F$ lies above $j$, and $j<n$ because $n\in F$.
There are $L=n-j-1$ slots below the maximum. Their digits are in
$\{0,1,2\}$, and their sum is exactly $K-1$: nothing is lost below the
cutoff, and precisely the unique maximum is omitted. Consequently,
\[
 c_{q,r}(E_{q,r,n}(F))
   =n-j-1+\floor{\frac{K-1+r}{q}}=n-1.
\]

Now let $w$ have cost $n-1$ and set $j=\floor{(S+r)/q}$. The cost
identity says $n=L+j+1$. The proposed lower labels $j+1,\ldots,j+L$
are therefore positive, distinct slots below $n$, each used at most twice.
Together with one $n$ they form a submultiset of $U_n$. Write
\[
 S+r=qj+u,\qquad 0\le u<q.
\]
The reconstructed multiset $F$ has cardinality $S+1$, and all its labels,
including $n$, are at least $j+1$. Hence
\[
 |F|+r=qj+u+1\le q(j+1)\le q\min F,
\]
so the reconstructed multiset belongs to $\A_{q,r}(n)$.

Starting with $F$, the word sum is $K-1$, so decoding recovers the same
cutoff $j$ and every multiplicity at every lower slot, followed by its one
maximum. Thus decoding after encoding is the identity. Starting with a
word, the decoded cardinality is $S+1$, so re-encoding again uses the same
$j$. Its maximum is $L+j+1$, so re-encoding reads exactly $L$ positions and
returns each digit $d_t$, including all zeros. This proves the other
inverse identity.
\end{proof}

The word family is finite independently of the multiset description:
$c_{q,r}(w)\ge L$, so a word of cost $n-1$ has length at most $n-1$ and
belongs to a set of size $1+3+\cdots+3^{n-1}$. In particular, all word counts used below are finite.

At $n=1$, the only multiset is $\{1\}$ and its word is empty, since
$1+r\le q$. For arbitrary $n$, the singleton $\{n\}$ instead gives the
zero word of length $n-1$. At $n=2$, a word of cost one must have exactly
one digit, which satisfies $d+r<q$. Therefore
\begin{equation}\label{eq:n2}
 a_{q,r}(2)=\min(3,q-r).
\end{equation}
These statements include $q=1$. Both initial and final zeros must be
retained. For example, in $\A_4(4)$ the multiset $\{2,2,4\}$ gives
$(0,2,0)$: its cutoff is zero, although its minimum is two. Removing
either zero changes the cost and the reconstructed multiset. The term
$K-1+r$ in the cutoff also matters when $K+r$ is divisible by $q$.

\section{Residue equations and their unique solution}

For a starting residue $r$ and a digit $d$, define
\begin{equation}\label{eq:carry}
 e=\floor{\frac{r+d}{q}},\qquad r'=(r+d)\bmod q.
\end{equation}
Write $dw$ and $wd$ for the words obtained by adjoining $d$ at the
beginning and at the end of $w$, respectively. Since $r+d=qe+r'$, direct
substitution in \eqref{eq:cost} gives
\begin{equation}\label{eq:digit}
 c_{q,r}(dw)=c_{q,r}(wd)
     =1+e+c_{q,r'}(w).
\end{equation}
The two words have the same cost, but need not be equal.

Put $C_{q,r}(N)=|\{w:c_{q,r}(w)=N\}|$ for $N\ge0$ and use
$C_{q,r}(N)=0$ for $N<0$.
The empty word is the unique word of cost zero. For a nonempty word
beginning with $d$, deleting that digit reduces the cost by
$1+\floor{(r+d)/q}$ and changes the starting residue to $(r+d)\bmod q$.
Prepending $d$ is the inverse operation. Partitioning by the first digit
therefore gives, for all $N\ge0$,
\begin{equation}\label{eq:coeff}
 C_{q,r}(N)=\mathbf1_{N=0}+
   \sum_{d=0}^2 C_{q,(r+d)\bmod q}
       \left(N-1-\floor{\frac{r+d}{q}}\right).
\end{equation}
The images for distinct first digits are disjoint, and terms with negative
indices vanish by the stated convention.

\begin{proposition}\label{prop:system}
For fixed $q\ge1$, the series
\[
 G_{q,r}(x)=\sum_{n\ge1}a_{q,r}(n)x^{n-1}
           =\sum_{N\ge0}C_{q,r}(N)x^N\quad(0\le r<q)
\]
are the unique solution in $\Z[[x]]^q$ of
\begin{equation}\label{eq:system}
 G_{q,r}=1+\sum_{d=0}^2
 x^{\,1+\floor{(r+d)/q}}G_{q,(r+d)\bmod q}.
\end{equation}
\end{proposition}

\begin{proof}
Theorem~\ref{thm:encoding} identifies the corresponding coefficients of
the two series. Equation~\eqref{eq:coeff} then yields the system
\eqref{eq:system}.
Every exponent multiplying a series on the right is at least one. Thus
all constant coefficients equal one, and at degree $N>0$ the right side
uses only coefficients of smaller degree. Induction on $N$ proves
uniqueness simultaneously for all residues. It also constructs a solution
coefficient by coefficient in the ring of formal power series.
\end{proof}

For $q\ge2$, the bounds $r<q$ and $d\le2$ imply $e\le1$. Translating
\eqref{eq:coeff} back to positive multiset indices gives, for $n\ge3$,
\begin{equation}\label{eq:positive}
 a_{q,r}(n)=\sum_{d=0}^2
 a_{q,(r+d)\bmod q}\left(n-1-\floor{\frac{r+d}{q}}\right).
\end{equation}
All three arguments are positive. For $q=1$ the digit $d=2$ instead gives
a carry of two; the word correspondence and \eqref{eq:coeff} still apply,
with the same convention for negative coefficient indices.

\section{Nine polynomial identities}

Let
\begin{align*}
 D_2&=1-x-2x^2-x^3,\\
 D_3&=1-3x+3x^2-4x^3+2x^4-x^5,\\
 D_4&=1-3x+3x^2-3x^3-2x^4-x^5.
\end{align*}
With $h=1-x+x^2$, define the numerator vectors, in residue order, by
\begin{align*}
 (N_{2,0},N_{2,1})&=(1+x,1),\\
 (N_{3,0},N_{3,1},N_{3,2})&=(1,h,h^2),\\
 (N_{4,0},N_{4,1},N_{4,2},N_{4,3})
   &=(1+2x^2,\ 1+x^3,\ 1-x+x^2+x^3+x^4,\ 1-2x+3x^2).
\end{align*}
These polynomials satisfy, for each residue,
\begin{equation}\label{eq:certificate}
 N_{q,r}-\sum_{d=0}^2x^{\,1+\floor{(r+d)/q}}
                 N_{q,(r+d)\bmod q}=D_q.
\end{equation}
For completeness, the nine rows are the following explicit identities.
For $q=2$,
\begin{align*}
 (1-x-x^2)(1+x)-x&=D_2,\\
 (1-x-x^2)-x^2(1+x)&=D_2.
\end{align*}
For $q=3$,
\begin{align*}
 (1-x)-xh-xh^2&=D_3,\\
 (1-x)h-xh^2-x^2&=D_3,\\
 (1-x)h^2-x^2-x^2h&=D_3.
\end{align*}
For $q=4$, abbreviating $N_{4,r}$ to $N_r$,
\begin{align*}
 (1-x)N_0-xN_1-xN_2&=D_4,\\
 (1-x)N_1-xN_2-xN_3&=D_4,\\
 (1-x)N_2-xN_3-x^2N_0&=D_4,\\
 (1-x)N_3-x^2N_0-x^2N_1&=D_4.
\end{align*}
Each identity follows by expanding the specified polynomials in $\Z[x]$.

\begin{proposition}\label{prop:gf}
For $q=2,3,4$ and $0\le r<q$, $G_{q,r}=N_{q,r}/D_q$ as formal power
series. In particular,
\begin{equation}\label{eq:gf}
 \sum_{n\ge1}a_2(n)x^{n-1}=\frac{1+x}{D_2},\qquad
 \sum_{n\ge1}a_3(n)x^{n-1}=\frac1{D_3},\qquad
 \sum_{n\ge1}a_4(n)x^{n-1}=\frac{1+2x^2}{D_4}.
\end{equation}
\end{proposition}

\begin{proof}
Each $D_q$ has constant coefficient one and hence has a unique inverse in
$\Z[[x]]$: its coefficients are obtained recursively from $D_qD_q^{-1}=1$.
Dividing \eqref{eq:certificate} by $D_q$ shows that the vector
$(N_{q,r}/D_q)_r$ satisfies \eqref{eq:system}. Uniqueness from
Proposition~\ref{prop:system} identifies it with the generating functions of the
multiset counts.
\end{proof}

\begin{proof}[Proof of Theorem~\ref{thm:main}]
Write $b_k=a_q(k+1)$, so that the series on the left of \eqref{eq:gf} is
$\sum_{k\ge0}b_kx^k$. If $D_q=1+\sum_{i=1}^m\delta_i x^i$, comparison
of the coefficient of $x^k$ in $D_qG_{q,0}=N_{q,0}$ gives
\begin{equation}\label{eq:extract}
 b_k+\sum_{i=1}^{\min(k,m)}\delta_i b_{k-i}=[x^k]N_{q,0}.
\end{equation}
For $q=2$, the first three equations give $b_0=1$, $b_1-1=1$, and
$b_2-2-2=0$, hence $1,2,4$. For $q=3$, the corresponding equations give
\[
 b_0=1,\quad b_1=3,\quad b_2=9-3=6,\quad
 b_3=18-9+4=13,\quad b_4=39-18+12-2=31.
\]
For $q=4$, the numerator coefficient two at degree two gives
\[
 b_0=1,\quad b_1=3,\quad b_2=9-3+2=8,\quad
 b_3=24-9+3=18,\quad b_4=54-24+9+2=41.
\]
For $k\ge3$ in the first case, and $k\ge5$ in the other two cases, the
right side of \eqref{eq:extract} vanishes and all $m$ earlier coefficients
occur. Moving those terms to the right and putting $n=k+1$ gives exactly
\eqref{eq:r2}--\eqref{eq:r4}, starting at $n=4,6,6$, respectively.
\end{proof}

\begin{corollary}\label{cor:minimal}
The minimal orders of eventual homogeneous constant-coefficient recurrences
over $\Q$ for $a_2,a_3,a_4$ are $3,5,5$, respectively.
\end{corollary}

\begin{proof}
The fractions in \eqref{eq:gf} are reduced over $\Q[x]$. For $q=2$,
$D_2(-1)=1$; for $q=3$, the numerator is one. For $q=4$, polynomial
subtraction gives
\[
 4D_4=(1+2x^2)(8-5x-4x^2-2x^3)-4-7x.
\]
A nonconstant common divisor would divide $4+7x$, but
$1+2(-4/7)^2=81/49\ne0$, so there is none.
Now an eventual recurrence of order $d$ gives a nonzero polynomial $R$
of degree at most $d$, with constant coefficient one, such that
$RG_{q,0}$ is a polynomial: every sufficiently high coefficient vanishes.
Writing $G_{q,0}=N/D$ in reduced form, $D\mid RN$ and coprimality imply
$D\mid R$. Therefore $d\ge\deg D$, regardless of where the recurrence
begins. Theorem~\ref{thm:main} attains these three bounds.
\end{proof}

\section{The specialised constructions}

\subsection{Compositions with a distinguished terminal part}

A composition is a nonempty ordered list of positive integers with a
specified total. The terminal part in the next statement is constrained
differently from the other parts; this distinction explains the special
simplicity of $q=3$.

\begin{proposition}\label{prop:terminal}
For $q\ge1$, $0\le r<q$, and $n\ge1$, the family $\A_{q,r}(n)$ is in
bijection with the compositions of $qn+q-1-r$ whose nonterminal parts
belong to $\{q,q+1,q+2\}$ and whose terminal part belongs to
$\{q,q+1,\ldots,2q-1\}$.
\end{proposition}

\begin{proof}
First apply Theorem~\ref{thm:encoding}. For its word $w=(d_1,\ldots,d_L)$,
put $S=S(w)$ and $u=(S+r)\bmod q$. Send it to
\begin{equation}\label{eq:composition}
 (q+d_1,\ldots,q+d_L,\ 2q-1-u).
\end{equation}
The final part is in the stated range. With $S+r=qj+u$ and
$L+j=n-1$, its total is
$qL+S+2q-1-u=qn+q-1-r$.

To prove surjectivity and the inverse identity, take any composition of
that total satisfying the two part restrictions. Let $\ell$ be its
length and $t$ its final part. All parts are at least $q$, and the total
is less than $q(n+1)$, so $1\le\ell\le n$. Subtract $q$ from the
$\ell-1$ nonterminal parts to obtain a ternary word $w$. The total
identity rearranges to
\[
 S(w)+r=q(n-\ell)+(2q-1-t).
\]
Here $0\le2q-1-t<q$. Thus $\floor{(S(w)+r)/q}=n-\ell$, and the word
has cost $(\ell-1)+(n-\ell)=n-1$. Its remainder is $2q-1-t$, so
\eqref{eq:composition} recovers the given terminal part as well as every
nonterminal part. Conversely, subtracting $q$ after the forward map
recovers each original digit. These maps are therefore inverse bijections
between words of cost $n-1$ and the stated compositions. Composing with
the bijection in Theorem~\ref{thm:encoding} proves the result.
\end{proof}

\begin{corollary}\label{cor:q3}
For every $n\ge1$, $\A_3(n)$ is in bijection with all compositions of
$3n+2$ into parts $3,4,5$.
\end{corollary}

\begin{proof}
For $q=3,r=0$, the terminal and nonterminal alphabets in
Proposition~\ref{prop:terminal} are both $\{3,4,5\}$.
\end{proof}
For a multiset of cardinality $K$, the final part is exactly
$5-((K-1)\bmod3)$. At $n=2$, the correspondence is
\[
 \{2\}\longleftrightarrow(3,5),\qquad
 \{1,2\}\longleftrightarrow(4,4),\qquad
 \{1,1,2\}\longleftrightarrow(5,3).
\]
Thus $a_3(n)$ is the term of A017818 at index $3n+2$ \cite{OEIS017818}.
For other values of $q$, the two part restrictions need not coincide: for $q=2$ the
terminal alphabet is $\{2,3\}$ rather than $\{2,3,4\}$; for $q=4$ it is
$\{4,5,6,7\}$ rather than $\{4,5,6\}$. Already $q=4,n=1$ gives the
composition $(7)$.

\subsection{The six maps for the strict companion}

Put $B_n=\A_{2,1}(n)$, so its members satisfy $|F|<2\min F$, and write
$A_n=\A_2(n)$. We now realise \eqref{eq:positive} by explicit multiset maps, using
the last digit of the word.

For $q\ge2$, $n\ge3$, and a fixed target residue $r$ and last digit $d$,
let $e,r'$ be as in \eqref{eq:carry}, and put $n'=n-1-e$. From a source
$F\in\A_{q,r'}(n')$, define
\begin{equation}\label{eq:transport}
 T_{r,d}(F)=\sigma^e(F\setminus\{n'\})
                \uplus d\{n-1\}\uplus\{n\},
\end{equation}
where $\sigma^e$ adds $e$ to every label, the deletion removes one
occurrence, and $d\{n-1\}$ denotes $d$ copies. The source maximum is
unique, so it is completely removed before shifting.

\begin{proposition}\label{prop:maps}
The map \eqref{eq:transport} is a bijection onto the part of
$\A_{q,r}(n)$ with multiplicity $d$ at $n-1$. Under the word
correspondences it is $w\mapsto wd$.
\end{proposition}

\begin{proof}
Let $K=|F|$ and $j'=\floor{(K-1+r')/q}$. The target has cardinality
$K+d$, and its cutoff computed by the encoding formula is
\[
 \floor{\frac{K+d-1+r}{q}}
   =e+\floor{\frac{K-1+r'}q}=e+j'.
\]
Every lower source slot shifts by $e$ and retains its digit. Its last
possible label is $n'-1+e=n-2$, so the new slot at $n-1$ has exactly
$d$ occurrences, followed by the unique maximum $n$. Equivalently,
\eqref{eq:transport} is the decoding of the source word with $d$
appended. Equation~\eqref{eq:digit} makes its cost $n-1$, so the decoder
proves admissibility as well.

Conversely, a target of cost $n-1>0$ has a nonempty word; its last slot is
at $n-1$, so its last digit is exactly that multiplicity. Removing $d$
leaves a word of cost $n-2-e=n'-1$. Theorem~\ref{thm:encoding} decodes
it to a unique member of the specified source family. Appending and
removing this digit are inverse operations, proving the assertion.
\end{proof}

For $q=2$, Proposition~\ref{prop:maps} gives the six restricted maps in
one table. Its entries are source families, for each prescribed target
multiplicity at $n-1$:
\[
\begin{array}{c|ccc}
\text{target}&d=0&d=1&d=2\\ \hline
 A_n&A_{n-1}&B_{n-1}&A_{n-2}\\
 B_n&B_{n-1}&A_{n-2}&B_{n-2}.
\end{array}
\]
For $d=0$ the operation replaces the old maximum by $n$. In the
one-copy relaxed case it is $F\mapsto F\uplus\{n\}$; in the one-copy
strict case it is $F\mapsto\sigma F\uplus\{n\}$. Both two-copy cases
are $F\mapsto\sigma F\uplus\{n-1\}\uplus\{n\}$. These equal
\eqref{eq:transport} because the shifted source maximum contributes the
first copy of $n-1$. In particular, writing $b(n)=|B_n|$, the two
companion equations are, for $n\ge3$,
\begin{align*}
 a_2(n)&=a_2(n-1)+b(n-1)+a_2(n-2),\\
 b(n)&=b(n-1)+a_2(n-2)+b(n-2).
\end{align*}

Related coloured families occur in \cite{CIMSZ}. Their distinction from
the present multisets is already visible at small indices. At maximum three, the two-coloured family satisfying
$\min C\ge|C|$ has members $\{3\},\{2_1,3\},\{2_2,3\}$, whereas
$B_3$ has members $\{3\},\{2,3\},\{2,2,3\}$. Forgetting colours
identifies two members and misses the double occurrence. Thus equality of these counts does not arise by simply forgetting colours.

For $q=4$, the four residue series count the four margin families
$\A_{4,r}(n)$. Partitioning their words by the first digit gives
\eqref{eq:system}; adjoining a last digit gives the maps
\eqref{eq:transport}. These constructions have the same cost equations,
but in general give different words and different multisets.

\section*{Formal verification}

Formal proofs of Theorems~\ref{thm:main} and~\ref{thm:encoding} accompany
this paper in Lean~4 \cite{Lean4}, using the mathematical library mathlib
\cite{Mathlib}. The development also verifies the general terminal-composition
bijection, the finite residue partition, and the correspondence with
the three case constructions. The formal-power-series development
constructs the actual counting series, proves uniqueness of the residue
solution and all nine rational identities, and verifies coefficient
extraction at every index. It also proves the exact minimal eventual
recurrence orders over the rationals, allowing arbitrary starting
indices. The supplementary material
contains the source code, a correspondence between the mathematical
statements and Lean declarations, and instructions for reproducing
the formal and computational checks.

\paragraph{Use of artificial intelligence tools.}
OpenAI ChatGPT and Codex were used in developing the proofs, preparing
the Lean formalisations, and drafting and editing the manuscript.

\bibliographystyle{amsplain}
\bibliography{references}
\end{document}
